\documentclass[11pt]{article}
\usepackage{microtype,amsmath,amssymb,amsthm,mathtools,booktabs,array}
\usepackage[margin=1.05in]{geometry}
\usepackage[colorlinks=true,linkcolor=blue!45!black,citecolor=blue!45!black,urlcolor=blue!45!black]{hyperref}
\usepackage{xcolor}
\hypersetup{pdftitle={The dimension-degree boundary for Gaussian moment vanishing},pdfsubject={Working mathematical manuscript}}
\newtheorem{theorem}{Theorem}[section]
\newtheorem{proposition}[theorem]{Proposition}
\newtheorem{lemma}[theorem]{Lemma}
\newtheorem{corollary}[theorem]{Corollary}
\theoremstyle{definition}\newtheorem{definition}[theorem]{Definition}
\theoremstyle{remark}\newtheorem{remark}[theorem]{Remark}
\newcommand{\E}{\mathbb E}
\newcommand{\C}{\mathbb C}
\newcommand{\tr}{\operatorname{tr}}
\newcommand{\wt}{\operatorname{wt}}
\newcommand{\G}{\mathsf G}
\DeclareMathOperator{\Span}{span}
\title{A resolved-curve logarithmic seed bound for Laurent phases}
\author{}\date{17 November 2026: retained written proof}
\begin{document}\maketitle
This proof uses the convergent nearby-cycle and period-comparison
theorems credited below. It supplies a complete marked model,
nearby-cycle chain construction and logarithmic and tail bounds for
the distinguished differential $d\log x$. It is a written internal
argument,without independent replay,machine-formal verification,
external review or novelty claim. It does not concern arbitrary
nonconvergent Laurent families or all-ring descent.
\section{Gaussian conventions for the applications}
Let $T$ be standard Gaussian and $Z,W$ circular coordinates with
$\E[Z^iW^j]=\delta_{ij}i!$, independent of $T$. For a polynomial
$P$, set $\Phi_Q=\sum_{m\geq0}\E[QP^m]t^m/m!$. Purity means
$\Phi_1=1$; safety means every $\Phi_Q$ is a polynomial.
These are formal series, without an integral-convergence hypothesis.
Gaussian integration by parts is $\E[(\partial_T-T)H]=0$.

For any polynomial $h(z,T)$, let $g=th(g,T)$ be the unique small
formal solution. The exact circular identity is
\[
 \E_{Z,W}[A(Z,T)e^{tWh(Z,T)}]
       =\frac{A(g,T)}{1-th_Z(g,T)}.
\]
Indeed the left side is $\sum_m t^m[z^m]A(z,T)h(z,T)^m$.
It is the formal residue of $A(z,T)/(z-th(z,T))$, whose unique
small pole is $g$. This also proves the identity for a
$t$-dependent polynomial-coefficient formal numerator $A$.


\section{Seed restriction and curve estimate}
\begin{lemma}[The seed functional]\label{lem:seedfunctional}
Let $M$ be a finite-dimensional differential module over $K=\C(t)$,
let $s\in M$, and suppose a $K$-linear horizontal map
$J:M\longrightarrow\C((t))$ satisfies $J(s)=1$. Then the differential
submodule $U=K\langle\partial_t\rangle s$ has a horizontal quotient
$L:U\longrightarrow K$ with $L(s)=1$.

At an ordinary base point, $L$ extends locally to an analytic horizontal
functional on $M$. For the period conclusion, assume explicitly that
$M$ is the generic exponential de Rham system of regular functions
$f,g$ on a smooth affine complex variety, where $f$ is the base
projection and $f,g$ are algebraically independent. Remove the finite
exceptional base set in the Hien--Roucairol comparison theorem.
Then the extension is represented by a flat rapid-decay cycle, and its
period evaluated on $s$ is identically one.
\end{lemma}
\begin{proof}
On the generators $s,\nabla s,\nabla^2s,\ldots$, the values of $J$
are $1,0,0,\ldots$. Therefore its restriction takes values in $K$.
Every rational relation between these generators is respected because
$J$ was defined on $M$, so this restriction is well defined and
horizontal. Its value at $s$ makes it surjective.

Choose a point where $M$, $U$ and $L$ are ordinary, extend the fiber
functional from $U$ to $M$, and solve the dual analytic ODE.
Uniqueness preserves its restriction to $U$. The perfect period
pairing and solution representation in
Hien--Roucairol~\cite[Theorems~2.4 and~3.7]{HienRoucairol} give the
asserted cycle. A scalar rational normalization is an invertible gauge
away from its zeros and poles and must be retained in the period.
\end{proof}

\begin{lemma}\label{lem:curvelength}
A real algebraic curve of degree at most $D$ has length at most
$\pi D\rho$ in a disk of radius $\rho$. Consequently its logarithmic
length in $a\leq|x|\leq b$ is at most a constant times
$D(1+\log(b/a))$.
\end{lemma}
\begin{proof}
Average the number of intersections with lines over directions and
apply the one-dimensional coarea formula. Outside exceptional line
components the number is at most $D$; the components themselves
satisfy the same estimate. In each dyadic annulus
$\rho\leq|x|\leq2\rho$, divide the disk estimate by $\rho$ to
bound $\int|dx|/|x|$ by $2\pi D$. Sum the dyadic annuli.
\end{proof}


\subsection{Resolving a Laurent phase}\label{sec:laurentresolution}

The logarithmic estimate needs two ingredients: a complete description
of the regular part, and a bound for the distinguished differential
$d\log x$. The first is a nearby-cycle calculation. The second is a
calculation on annuli. Keeping these ingredients separate prevents a
perfect period pairing from being mistaken for a growth estimate.

Throughout this subsection the coefficients are rational functions of
the parameter. In particular their Puiseux expansions converge. Arbitrary nonconvergent formal Laurent families and arbitrary
coefficient rings are outside this proof scope.

The two-ended twisted quotient is finite dimensional without a
genericity assumption. If the actual Laurent orders are $p,q\geq1$,
the classes $x^{-q},\ldots,x^{p-1}$ are a basis. The top term of
$x\varphi_x$ eliminates every exponent at least $p$, without crossing
below $-q$; the bottom term eliminates every exponent below $-q$,
without crossing above $p-1$. An exact expression supported in this
interval would require its preimage to have highest exponent at most
$-1$ and lowest exponent at least zero. This proves independence
and rank $p+q$. Coefficient specializations use their actual orders.

\begin{lemma}[A marked model]\label{lem:markedmodel}
Let $\varphi(r,x)\in\C(r)[x,x^{-1}]$ be nonconstant in $x$.
After a finite ramification of $r$ there is a smooth surface proper
over a parameter disc, unchanged over its punctured disc, with these
properties:
\begin{enumerate}
\item Its generic fibre is $\mathbb P^1$. Its central fibre is a
reduced normal-crossing curve, with local equations $r=u$ or $r=uv$.
\item The phase extends to a morphism to $\mathbb P^1$. Its zero
and pole divisor and the divisor of $x$ have normal crossings with
the central fibre.
\item The horizontal sections $x=0,\infty$ are disjoint and meet
smooth points of the central fibre. Locally $x=u^a v^b U$, where
$a,b$ are integers and $U$ is holomorphic and nonvanishing.
\end{enumerate}
Additional markings used in the construction are not additional
punctures of the exponential connection.
\end{lemma}
\begin{proof}
Mark $0,\infty$ and every point of the generic fibres
$\varphi^{-1}(0),\varphi^{-1}(1),\varphi^{-1}(\infty)$.
Retain multiplicities in the divisors, but identify identical sections.
Convergent Newton--Puiseux makes the finitely many marked branches
meromorphic sections after one finite ramification. Start with
$\mathbb P^1$ over that disc and blow up collisions of distinct
sections. This terminates: the contact order of each pair is a finite
integer and decreases on blowing up their common point.

Every centre is a smooth point of a single central component. Indeed
a section cannot pass through a node $uv=r$: two positive integral
orders cannot have sum one. Thus each new exceptional component has
multiplicity one and all nodes still have the equation $uv=r$.
The final marked sections are disjoint and miss the nodes.

The divisor of the meromorphic function $\varphi$ is supported on
the central components and its marked zero and pole sections. If a
point were an indeterminacy, its local divisor would have a zero
branch and a pole branch meeting there. After cancelling common
factors its expression is a unit times $u^a/v^b$, $a,b>0$.
The equation $\varphi=1$ then has a horizontal branch through that
point. At a node this contradicts the section property; at a smooth
point it contradicts the separation from the zero or pole section.
These are all possible intersections. Hence $\varphi$ is a morphism.
The asserted normal crossings and the expression for $x$ follow
from its divisor: horizontally $x$ has just the two marked sections.
Only those original sections are removed from the generic fibre;
marking a zero or a point with phase value one does not remove it.
\end{proof}

The marking of phase value one is necessary. For $\varphi=r/x^2$,
the zero and pole sections are already disjoint, but the map is
undefined at $r=x=0$. Its value-one branches $x=\pm\sqrt r$
force the cover $r=s^2$ and the chart $x=sy$, where the phase is
$1/y^2$. The marked construction therefore detects a transition
that merely separating zeros and poles would miss.


We use the following precise comparison facts. For a meromorphic
connection on a punctured disc, its moderate nearby cycles describe
its formal regular summand, with monodromy. Nearby cycles commute
with proper direct image. For a rank-one exponential on a surface
with normal crossings, a component on which the phase has a vertical
pole contributes no moderate nearby cycles. At a node between a
finite-phase component and a pole component, the surviving object is
the exponential connection on the finite component, localized at the
node; it is not an extra point class. At a node between two finite
components, the local complex is the usual nearby annulus complex,
including its monodromy. These are standard comparison and
normal-crossing facts, not new theorems here; see
Roucairol~\cite[Section~1.2 and Lemma~1.4]{RoucairolFormal} and
Sabbah~\cite[Section~1.b, especially Proposition~1.5]{SabbahStationary}.

For clarity, their local exponential models are a holomorphic unit
(finite phase), $\exp(u^{-p})$ at a finite--pole boundary, and
$\exp(u^{-p}v^{-q})$, $p,q>0$, when both components are poles.
An analytic unit in a pole monomial is absorbed in a coordinate.
This may multiply the equation of the base by a nonvanishing unit;
locally that change does not change the vanishing assertion or the
nearby-cycle object up to the usual unit identification. In the
$f^s$ description the identification is multiplication by
$\exp(s\log U)$, with a local analytic choice of $\log U$.
One keeps the actual gluing maps and monodromy, rather than just
adding the ranks of the components.

\begin{lemma}[Lifting the regular part]\label{lem:regularlifting}
For the marked model, the complete nearby-cycle space of the
exponential direct image is represented, on a sufficiently small
parameter sector, by flat rapid-decay chains made from the following
pieces:
\begin{enumerate}
\item bounded deformations of finitely many paths and cycles in
the finite-phase central components, away from their nodes and poles;
\item a bounded number of bridges and circles in the necks
$uv=r$ between finite-phase components;
\item descending phase tails at irregular ends and circles about
regular punctures.
\end{enumerate}
The coefficients of each chosen flat chain are fixed constants.
Their nearby-cycle classes exhaust the formal regular part, with
its actual gluing maps and monodromy.
\end{lemma}
\begin{proof}
Apply proper nearby-cycle comparison to the compactified surface,
with the exponential connection localized at $r=0$ and at the two
original punctures. The normal-crossing rules just stated kill the
interiors of the pole components. Along a finite component they leave
its ordinary exponential connection, with localization at its pole
ends and at any original regular puncture. A finite--pole node gives
that localized end. A finite--finite node retains the nearby annulus.
Thus no pole component is silently counted as an additional regular
class, and no annulus or puncture residue is discarded.

Here is a chain description of this calculation. Remove small discs
about the finitely many nodes and punctures in each finite component.
Triangulate the remaining compact surfaces, keeping the boundary
circles. At an irregular boundary attach the finitely many descending
sectors. At a regular puncture keep the circle. The one-dimensional
exponential period comparison identifies the resulting relative
chain complexes with the component de Rham complexes; for a constant
phase this is ordinary homology. At each finite--finite node glue by
the actual annulus chain complex, with both its degree-zero and
degree-one terms. Taking the total complex is Mayer--Vietoris for
the full nearby-cycle complex above. It computes its maps and
monodromy, not merely its Euler characteristic. In particular its
classes exhaust the formal regular summand by proper comparison.

Choose a finite basis of this finite-dimensional homology. Smooth
pieces and their boundary points lift by the implicit function
theorem; an annulus is lifted by $uv=r$. Descending ends are continued
by the inverse phase coordinate. On a fixed small sector these choices
are isotopies, so the homology coefficients are constants. No Morse
condition on the phase is used. Degenerate critical points stay
inside the triangulated pieces, and a constant phase on a component
is allowed.

For the tail continuation, choose the ray of $r$ so that no unbounded
critical value lies on the negative real phase axis for small $r$.
There are finitely many excluded directions: critical values are
convergent Puiseux functions. Bounded critical values fit in a fixed
disc. From a sufficiently negative phase value every chosen inverse
branch therefore continues to $-\infty$ without meeting a critical
value. In the two-ended Laurent case it ends at $x=0$ or $\infty$.
The same choice can avoid the finitely many limiting Stokes directions
of the parameter connection.

This constructs flat lifts whose nearby-cycle classes exhaust the
space computed by proper comparison. No quantitative estimate has
been used in that completeness calculation.
\end{proof}

\begin{lemma}[The logarithmic bound on the lifted pieces]
\label{lem:logpieces}
For the chains of Lemma~\ref{lem:regularlifting}, the period of
$e^{\varphi}d\log x$ is $O(1+|\log r|)$ along the chosen ray.
The bound is uniform for a fixed finite collection of generators.
With a fixed Laurent amplitude the same construction has at most
an additional algebraic power of $|r|^{-1}$.
\end{lemma}
\begin{proof}
On a compact finite-phase piece the phase and the coefficients of
$d\log x$ are bounded, so its integral is bounded. In a finite--finite
neck write
\[
 r=uv,\qquad x=u^a v^b U(u,v).
\]
At fixed $r$,
\[
 d\log x=\left(a-b+\frac{uU_u-vU_v}{U}\right)\frac{du}{u}.
\]
The coefficient is uniformly bounded in a small bidisc. A radial
bridge across $|r|/\eta\leq|u|\leq\eta$, with bounded angular
detours, has logarithmic length $O(1+|\log r|)$. Circles have bounded
logarithmic length. The phase is holomorphic on this bidisc and is
uniformly bounded there. This treats every finite--finite neck.

Near an irregular end of a finite component, use a local root of
the phase unit to write $\varphi=w^{-p}$, $p\geq1$. At a
finite--pole node one can simultaneously write $r=wv$ after changing
the other coordinate by a unit. On $|v|,|w|<\eta$, the same
calculation bounds $d\log x$ by a constant times $|dw|/|w|$.
Along a descending phase ray, $s=-\varphi>0$, its integral is bounded
by a constant times
\[
 \int_{s_0}^{\infty}e^{-s}\frac{ds}{s}.
\]
If the node chart ends, truncate it first at $|w|=|r|^\alpha$,
where $0<\alpha<1$. Its phase is already of size
$-|r|^{-\alpha p}$. Continue that inverse branch along the negative
phase axis as in Lemma~\ref{lem:regularlifting}.

We give a uniform bound for that continuation, rather than declaring
it negligible from a local saddle formula. Choose $M$ large enough
to include all coefficient transition scales in
$|r|^M\leq|x|\leq|r|^{-M}$. Such $M$ exists because the coefficients
are rational and there are finitely many Laurent powers. In this
annulus each descending path lies on $\operatorname{Im}\varphi=0$,
a real algebraic curve of bounded degree after clearing its Laurent
denominator. The curve-length bound used in
Lemma~\ref{lem:curvelength} gives
$\int|dx|/|x|=O(1+|\log r|)$, by dyadic annuli. Along the entire
tail the exponential is at most
$\exp(-|r|^{-\alpha p})$. Outside that annulus the leading Laurent
monomial dominates uniformly, so
$|d\log x/d\varphi|\leq C/|\varphi|$ and its descending integral
is bounded by the same exponential tail integral. This proves the
required uniform tail bound. A fixed Laurent amplitude contributes
only a power of $|r|^{-1}$ on the annulus and algebraic growth in
$s$ outside; the exponential still dominates it.

At a regular puncture use a small fixed circle in its smooth marked
chart. Both its phase and logarithmic length are bounded. There are
finitely many pieces and their chain coefficients are fixed.
Adding their bounds proves the lemma.
\end{proof}

\begin{theorem}[Logarithmic seed bound]\label{thm:laurentlog}
Let $\varphi(t,x)\in\C(t)[x,x^{-1}]$. On a generic ray at
$t=\infty$, the zero-exponential part of every scalar flat period
of $e^{\varphi(t,x)}dx/x$ is $O(1+\log|t|)$. This includes all
critical degeneracies, coincident actions and regular punctures.
Consequently, for a rational normalization $h(t)$ with
$v_\infty(h)>0$, its differential seed submodule admits no rational
horizontal quotient taking the seed to one. The same exclusion
holds for a horizontal formal Gaussian functional at another point.
\end{theorem}
\begin{proof}
If the phase is independent of $x$, the sole residue period is a
constant times $e^{\varphi(t)}$. Its zero-exponential part is bounded
when $\varphi$ is regular at infinity, and absent otherwise.
For a nonconstant phase put $r=1/t$ and use
Lemmas~\ref{lem:markedmodel}--\ref{lem:logpieces}. The construction
includes regular puncture circles when there is just one irregular end.
Finite ramification changes $\log|r|$ only by a constant factor.
The lifted chains have at most algebraic growth on every Laurent
amplitude by Lemma~\ref{lem:logpieces}. On the chosen generic ray,
a full period vector with such growth has no exponentially growing
factor: after formal exponential decomposition a largest growing
factor could not cancel on every basis amplitude. The other nonzero
factors are exponentially decaying there. The nearby-cycle basis
therefore gives a full basis of regular functionals modulo those
decaying terms, with constant coefficients in the flat basis. Every
one of their seed periods has the asserted logarithmic bound. Thus
their zero-exponential seed
expansions have no negative algebraic powers. Multiplication by $h$
makes all of them tend to zero.

For a constant phase the rank-one period already gives the same
contradiction directly. For a nonconstant phase, a quotient with
normalized seed one would extend locally to a
horizontal functional on the full module, by the seed restriction
and analytic extension argument of Lemma~\ref{lem:seedfunctional}.
Its comparison hypotheses hold on
$(\mathbb A^1\setminus S)\times\mathbb G_m$: remove the finitely
many parameter poles, zeros of $h$, and comparison exceptions.
The base projection and nonconstant phase are regular and
algebraically independent there.
Its scalar seed is constant. Distinct formal exponential factors
separate, so its zero-exponential seed is also one. This contradicts
the bound after normalization. A formal Gaussian functional with
unit seed first restricts to the same rational seed quotient, as in
the first part of Lemma~\ref{lem:seedfunctional}.
\end{proof}

The comparison theorems supply completeness here. The bounds come
from the logarithmic differential on the marked model and the explicit
descending-tail estimate. Neither part follows from a finite symbolic
calculation. This proof uses convergence of the rational parameter
family; it makes no assertion for arbitrary nonconvergent Laurent
families or arbitrary coefficient rings.

\section{The Gaussian application}
Let $h(Z,T)=a+bZ+cT+dZ^2+eZT+fT^2$,
$\Delta=e^2-4df\ne0$ and $P=Wh(Z,T)+G(Z,T)$,with arbitrary
polynomial $G$. All coordinate inverses below are checked on the
actual formal Gaussian branch.
\subsection{The logarithmic branch, including its removed component}

For $\Delta\ne0$, $f\ne0$, choose $\eta^2=\Delta$ and set
\[
 \rho=\frac{e+\eta}{2f},\quad \sigma=\frac{e-\eta}{2f},\quad
 u=\frac{\rho c-b+1/t}{\eta},\quad
 v=\frac{b-1/t-\sigma c}{\eta},\quad k=uv-a/f.
\]
On the circular constraint $z=th(z,T)$, put
$x=T+\rho z+u$, so that $T+\sigma z+v=k/x$. Then
\[
 z=\frac f\eta(x-k/x-u+v),\quad T=x-u-\rho z,
 \qquad
 \frac{dT}{1-th_Z}=\frac{dx}{t\eta x}.
\]
The coefficient of $t^{-2}$ in $k$ is $-1/\eta^2$, so the conic
is generically smooth. The exponent $tG(z,T)-T^2/2$ is Laurent
polynomial in $x$. On the formal Gaussian branch,
$x=1/(\eta t)+O(1)$, which makes every Laurent power legitimate.

If $f=0$, then $e\ne0$. Put $z_0=-c/e$ and $x=z-z_0$. With
$A(z)=a+bz+dz^2$, the constraint gives
\[
 T=\frac{z/t-A(z)}{ex},\qquad
 \frac{dT}{1-th_Z}=\frac{dx}{tex}.
\]
This localization is valid if $c\ne0$, or if $c=0$ and $a\ne0$.
If $a=c=0$, the chosen fixed point is instead $g=0$ identically;
inverting $g$ would discard the very branch being tested. In that
case the exact unit series is
\[
 \E_T\left[\frac{e^{tG(0,T)}}{1-t(b+eT)}\right].
\]
Use $x=T-(1/t-b)/e$, now a valid Laurent unit. Its volume is
$-dx/(tex)$ and its exponent is polynomial in $x$. These formulas
cover both the smooth conic and the removed-component branch.
We next connect these coordinates to the formal Gaussian functional.
This will close every one-ended phase without an analytic estimate.

\begin{lemma}[Formal logarithmic transport]\label{lem:logtransport}
In each of the three charts above, write $X=X(t,T)$ for the formal
coordinate, $D$ for the displayed circular denominator, and $B$ for
$tG(g,T)$ or $tG(0,T)$ on the removed component. Put
$\kappa=\eta,e,-e$, respectively. Then
\[
 J(q)=\E_T\!\left[\frac{q(X)e^B}{D}\right]
\]
is horizontal on the logarithmic twisted quotient with phase
$\varphi$ and normalization $h=1/t$, and $J(1)=\Phi_1$.
\end{lemma}
\begin{proof}
All Laurent powers of $X$ are legitimate by the unit checks above.
The coordinate identities give
\[
 X_T=t\kappa X/D,\qquad
 \varphi(t,X)=B-T^2/2.
\]
Consequently Gaussian integration by parts gives
$J((x\partial_x+x\varphi_x)q)=0$: the corresponding integrand,
multiplied by $t\kappa$, is
$(\partial_T-T)\{q(X)e^B\}$.

For parameter horizontality set $w=1/D$ and $R=X_t/X_T$.
Differentiating $w=X_T/(t\kappa X)$ gives the exact identity
\[
 w_t+w/t=\partial_T(wR).
\]
Subtracting $J(q_t+(\varphi_t-1/t)q)$ from $\partial_tJ(q)$ now
gives the Gaussian expectation of
$(\partial_T-T)\{wR q(X)e^B\}$, hence zero. These are formal
coefficientwise identities, without an integral-convergence hypothesis.
The seed equality is the circular contraction formula.
\end{proof}

\begin{proposition}\label{prop:nondeglog}
If $e^2-4df\ne0$, no polynomial $G$, of any degree, makes
$P=Wh(Z,T)+G(Z,T)$ pure.
\end{proposition}
\begin{proof}
The smooth-conic, linear-conic and removed-component charts are
exhaustive. Their Laurent-unit checks and
Lemma~\ref{lem:logtransport} give the formal Gaussian functional
with seed value $\Phi_1$ and normalization $1/t$.
Theorem~\ref{thm:laurentlog} excludes $\Phi_1=1$ in every chart,
using its actual Laurent orders. No image-kernel calculation is used.
\end{proof}


\begin{thebibliography}{99}
\bibitem{DVEZ}
H.~Derksen, A.~van den Essen and W.~Zhao,
\emph{The Gaussian Moments Conjecture and the Jacobian Conjecture},
\emph{Israel Journal of Mathematics} \textbf{219} (2017), no.~2,
917--928; \href{https://arxiv.org/abs/1506.05192}{arXiv:1506.05192}.

\bibitem{DLMF}
NIST Digital Library of Mathematical Functions,
\href{https://dlmf.nist.gov/10.25.E1}{10.25.1} and
\href{https://dlmf.nist.gov/10.31.E2}{10.31.2},
modified Bessel equation and logarithmic power series.

\bibitem{HienRoucairol}
M.~Hien and C.~Roucairol, \emph{Integral representations for solutions
of exponential Gau\ss--Manin systems},
\emph{Bulletin de la Soci\'et\'e Math\'ematique de France}
\textbf{136} (2008), no.~4, 505--532,
\href{https://doi.org/10.24033/bsmf.2564}{doi:10.24033/bsmf.2564}.

\bibitem{Liu}
Z.~Liu, \emph{Sharp finite Gaussian moments for quadratic polynomials},
version~3, public preprint (2026),
\href{https://doi.org/10.5281/zenodo.21688251}{doi:10.5281/zenodo.21688251}.

\bibitem{Long}
C.~D.~Long, \emph{Small Counterexamples to the Gaussian Moments Conjecture},
public preprint (2026),
\href{https://arxiv.org/abs/2607.18186v1}{arXiv:2607.18186v1}.

\bibitem{RoucairolFormal}
C.~Roucairol, \emph{Formal structure of direct image of holonomic
$\mathcal D$-modules of exponential type},
\emph{Manuscripta Mathematica} \textbf{124} (2007), no.~3, 299--318;
\href{https://arxiv.org/abs/math/0604134v2}{arXiv:math/0604134v2}.

\bibitem{SabbahStationary}
C.~Sabbah, \emph{An explicit stationary phase formula for the local
formal Fourier--Laplace transform}, in \emph{Singularities I},
Contemporary Mathematics \textbf{474}, American Mathematical Society,
2008, 309--330;
\href{https://arxiv.org/abs/0706.3570v2}{arXiv:0706.3570v2}.

\bibitem{Wall}
C.~T.~C.~Wall, \emph{Singular Points of Plane Curves}, Cambridge
University Press, 2004, Chapter~2, 15--38,
\href{https://doi.org/10.1017/CBO9780511617560.003}{doi:10.1017/CBO9780511617560.003}.

\bibitem{Wilson}
M.~Wilson, \emph{A face-isolation proof of the two-variable Gaussian
Moments Conjecture}, public preprint (2026),
\href{https://arxiv.org/abs/2607.23887v1}{arXiv:2607.23887v1}.
\end{thebibliography}
\end{document}
